% MT044 / explanation — canonical editable content.
随机森林训练许多有随机差异的树，再汇总预测。树之间的差异既来自抽样记录，也常来自每个节点只从部分候选特征中找切分点。



\begin{itemize}

\item \textbf{A按课堂含义正确。} 随机抽取候选特征能让各树不总依赖同一强特征，降低它们犯错的相关性。常见算法是在每个节点重新抽取，不一定给整棵树永久固定一组特征；原选项是简写。

\item \textbf{B错误。} 普通RF把各树预测平均或投票，没有boosting那种控制每轮新增模型贡献的学习率。树数、深度、候选特征数等才是常见RF超参数。

\item \textbf{C正确。} 分类树输出类别或概率，回归树输出数值，两类任务都可以构成森林；boosting也能针对分类或回归的损失训练。

\item \textbf{D错误。} 它把两个方法家族限制到单一任务，和C的事实相反。

\item \textbf{E错误。} A与C已成立。

\end{itemize}


\textbf{出处：} 2023A期中 Q5，原卷第3页（5分）；2025A期中 Q5，原卷第3页（5分）


\textbf{答案依据：} 2023A期中Q5随卷答案AC；2025A Q5同题。
